% =====================================================================
%  capitulos/cap14.tex — Capítulo 14. La epistemicidad: modalidad y
%  evidencialidad
%  Autora: Ana Marcovecchio
%  Fuente: Capitulo_14_normalizado.docx
% =====================================================================
\capitulo{La epistemicidad: modalidad y evidencialidad}{Ana Marcovecchio}

\section{Introducción}

En este capítulo se trata la modalidad, a partir de una definición gramatical
que opone el \textit{dictum} oracional, esto es, el contenido léxico de una
predicación, al \textit{modus}, o sea, la huella que imprime quien habla sobre
sus dichos, fundamentalmente, en relación con si los presenta, por ejemplo, como
aseveraciones, mandatos o deseos. La modalidad también comprende la polaridad,
ya que las predicaciones pueden ser afirmativas o negativas. Asimismo, la
manifestación de determinados posicionamientos modales está directamente
asociada al grado de conocimiento que quien enuncia tiene sobre los
acontecimientos a los que se refiere: esto se conecta con una noción que ha ido
ganando relevancia en la gramática, la de \termino{evidencialidad}. Si bien
\textit{modalidad epistémica} y \textit{evidencialidad} son nociones distintas,
interactúan en un dominio común que se conoce como el de la
\termino{epistemicidad}.

Al expresarse estos contenidos gramaticales mediante diferentes recursos, el
recorrido de este capítulo supone retomar varios de los temas ya vistos en este
libro: la morfología flexiva de los verbos; el potencial funcional de una clase
de palabras como el adverbio, que se caracteriza por la posibilidad de
comportarse como circunstancial de los verbos, pero también puede desempeñar
funciones extrapredicativas; las perífrasis verbales, sobre todo, las de
modalidad, y, en relación con esto, la organización oracional en tres capas (la
léxica, la flexiva y la informativa).

\section{Definición, delimitación y confluencia con otras categorías}\label{sec:14-definicion}

Indisociable de las propiedades de la estructura lingüística, aparece una
categoría cuyo deslinde acarrea controversia: la \termino{modalidad}. Esta
categoría lingüística codifica, mediante una notable variedad de recursos
gramaticales (véase §\ref{sec:14-recursos}), significados relativos a las
huellas que como hablantes imprimimos en nuestros enunciados, con respecto al
posicionamiento y la perspectiva que adoptamos sobre lo que decimos. Al
respecto, recordemos que, como se planteó en el \capref{6}, mientras que la
oración es una unidad de la gramática, el enunciado es su realización en
situaciones comunicativas concretas. En efecto, no existe la posibilidad de que
comuniquemos un determinado contenido sin que, en simultáneo, se articule con
una manera de enunciarlo: en esto consiste, en síntesis, la oposición entre el
\textit{dictum}, o contenido proposicional de la oración, y el \textit{modus}, o
modalidad, según la cual en todo enunciado (y en su forma abstracta, la oración)
se produce una combinación de un cierto contenido descriptivo con una
determinada intención de quien lo enuncia, que puede, por ejemplo, comunicar o
bien pedir información, pretender influir sobre su(s) interlocutor(es),
manifestar un deseo o formular una pregunta, entre otras
posibilidades.\pendiente{Comentario de la autora: no volver aquí sobre
«oración» frente a «enunciado»; señalar, en cambio, que la oración puede ganar
valores modales secundarios en el uso (p.~ej., interrogativas interpretadas como
pedidos)}

En su forma gramatical, las oraciones proyectan esta presencia de quien enuncia
en las tres capas estructurales que las conforman (véase el \capref{6},
§\ref{sec:06-capas}):

\begin{itemize}
  \item en la capa léxica, a través de la selección de las piezas que establecen
        la predicación sobre los acontecimientos, con su estructura argumental y
        la red de papeles temáticos que señalan a los participantes del evento;
  \item en la capa flexiva, mediante el anclaje temporal, que fundamentalmente
        indica simultaneidad, anterioridad o posterioridad del acontecimiento
        descripto por la predicación con respecto al \textit{hic et nunc} de la
        enunciación;
  \item en la capa complementante, de manera crucial, por ser esta la interfaz
        entre el contenido proposicional y su articulación con el discurso.
\end{itemize}

De ahí que podamos sostener que el \textit{dictum} oracional se estructura
también conforme al \textit{modus} que le imprime quien habla a través de
diferentes medios gramaticales. Pero antes de adentrarnos en ello, podemos
presentar las definiciones de modalidad que han propuesto
\textcite{kovacci1992} y \textcite{rodriguezramalle2005}, respectivamente.

\begin{cita}
La modalidad tiene carácter complejo: incluye al menos dos componentes: la
actitud adoptada con respecto al \textit{dictum} (declarativa, dubitativa, etc.)
y la indicación de la presencia del hablante como tal (como \sig{diciente}), que
implica su responsabilidad por la actitud \parencite[99]{kovacci1992}.
\end{cita}

\begin{cita}
La modalidad recoge las diferencias existentes entre oraciones como reflejo de
las distintas posturas que adopta el hablante, bien con respecto a la verdad del
contenido de la proposición que emite bien con respecto a la manera en que la
expresa. La modalidad se puede definir, por tanto, como la manifestación de las
actitudes u opiniones del hablante ante lo que dice y comprende
\parencite[470]{rodriguezramalle2005}.
\end{cita}

\subsection{Los recursos de la modalidad}\label{sec:14-recursos}

Los recursos que sirven para formalizar la modalidad oracional son múltiples y
variados:

\begin{itemize}
  \item La \textbf{entonación} o línea musical que, en la oralidad, se superpone
        al componente léxico-sintáctico organizado según la disposición de los
        diferentes ítems léxicos y gramaticales. En la escritura, la entonación
        queda esbozada en el empleo de los signos de entonación, y también de los
        de puntuación; así, la predicación de \exref{ex:14-juan1} se interpreta
        como una aserción; la de \exref{ex:14-juan2}, como una interrogación; la
        de \exref{ex:14-juan3}, como una reacción emotiva:
\end{itemize}

\begin{exe}
  \ex\label{ex:14-juan1} Viene Juan.
  \ex\label{ex:14-juan2} ¿Viene Juan?
  \ex\label{ex:14-juan3} ¡Viene Juan\ldots!
\end{exe}

\begin{itemize}
  \item La \textbf{selección del modo verbal}, expresada en la flexión del verbo,
        por ejemplo, en la oposición entre el modo indicativo y el modo
        subjuntivo, lo que conduce a interpretar aquello que comunica la
        profesora como una aserción en \exref{ex:14-profesora1} y como una orden
        en \exref{ex:14-profesora2}:
\end{itemize}

\begin{exe}
  \ex\label{ex:14-profesora1} La profesora dice que los estudiantes llegan
      temprano. (Aseveración).
  \ex\label{ex:14-profesora2} La profesora les dice a los estudiantes que lleguen
      temprano. (Pedido).
\end{exe}

\begin{itemize}
  \item El repertorio de \textbf{auxiliares de modalidad} presentes en las
        perífrasis verbales (véase el \capref{11}, §\ref{sec:11-modales}),
        unidades conformadas por dos verbos, uno ---en forma finita--- que
        aporta información de índole gramatical y otro ---en forma no finita---
        que manifiesta el contenido léxico, como expresiones de diferentes
        modalidades expresivas del hablante, \exref{ex:14-conducir} y
        \exref{ex:14-llover}:
\end{itemize}

\begin{exe}
  \ex\label{ex:14-conducir}
    \begin{xlist}
      \ex \textit{Puedo} conducir porque tengo licencia. (\sig{capacidad} o
          \sig{habilitación}).
      \ex \textit{Tengo que} conducir porque soy el conductor asignado.
          (\sig{necesidad}).
    \end{xlist}
  \ex\label{ex:14-llover} Mañana \textit{puede} llover. (\sig{posibilidad}).
\end{exe}

\begin{itemize}
  \item Ciertas \textbf{construcciones} también sugieren interpretaciones
        modales; así, las estructuras llamadas \textit{cuasirreflejas
        mediopasivas e impersonales} (véase el \capref{10},
        §§\ref{sec:10-mediopasivas} y \ref{sec:10-impersonales}), por ejemplo,
        se caracterizan por desencadenar lecturas deónticas
        \exref{ex:14-prenda} o epistémicas \exref{ex:14-incidentes}:
\end{itemize}

\begin{exe}
  \ex\label{ex:14-prenda} Esta prenda \textit{se lava a mano}. (\sig{debe ser
      lavada a mano}).
  \ex\label{ex:14-incidentes} A diferencia de los accidentes, los incidentes
      \textit{se evitan}. (\sig{son evitables} / \sig{pueden ser evitados}).
\end{exe}

\begin{itemize}
  \item En los márgenes de la predicación, la llamada \textbf{periferia}, que se
        configura en torno a la predicación primaria establecida a partir del
        verbo léxico, sus argumentos y sus circunstantes temporalizados, se
        caracteriza por alojar una serie de elementos que tienen ámbito
        habitualmente sobre todo el núcleo oracional y que expresan así diversos
        posicionamientos modales, materializados en sintagmas adverbiales,
        preposicionales o bien en oraciones subordinadas \exref{ex:14-nino}:
\end{itemize}

\begin{exe}
  \ex\label{ex:14-nino} \textit{Claramente} / \textit{por ahí} / \textit{según
      parece}, el niño descansó mal anoche: no deja de bostezar.
\end{exe}

Por caso, en \exref{ex:14-nino} se suscita la lectura de que los dichos sobre
cómo fue el descanso del niño se fundan, no en la comprobación directa de quien
habla, sino en una conclusión que obtiene a partir de ciertos indicios: los
continuos bostezos.

\begin{itemize}
  \item También se expresa la modalidad mediante \textbf{ítems léxicos
        específicos}, tales como nombres (\forma{posibilidad,
        probabilidad\ldots}), adjetivos (\forma{evitable, posible, probable,
        eventual\ldots}) y adverbios e índices de actitud (\forma{tal vez,
        acaso, posiblemente, capaz (que)\ldots}).
\end{itemize}

\section{La modalidad como expresión de la fuerza ilocutiva}\label{sec:14-ilocutiva}

De acuerdo con una mirada interlingüística, esto es, un enfoque en que se buscan
las propiedades comunes de las lenguas, es complejo definir la modalidad, dado
que las distintas lenguas no apelan exactamente a los mismos recursos ni
codifican el mismo tipo de contenidos modales. Sin embargo, lo que
\textcite{palmer1986} sugiere es que la distinción general entre contenido
proposicional y modalidad permite pensar en un paralelismo cercano al de
\termino{acto locutivo} y \termino{acto ilocutivo}. Esta última distinción,
propuesta por \textcite{austin1971}, se ha convertido en una herramienta
fundamental para estudiar la distancia entre lo que se dice literalmente y lo
que efectivamente se comunica, de modo que el acto locutivo supone la
articulación lingüística de un enunciado y la \termino{fuerza ilocutiva} refiere
a la intención o el propósito de ese acto de habla.

Por ejemplo, dos enunciados como \exref{ex:14-ordena} y \exref{ex:14-porque},
con distinta configuración gramatical (una oración exhortativa con verbo en modo
imperativo, en \exref{ex:14-ordena}; una interrogación negativa con \forma{por
qué} más verbo en modo indicativo, en \exref{ex:14-porque}), pueden servir a la
expresión de una misma intención: la de transmitir una orden, un pedido, una
sugerencia.

\begin{exe}
  \ex\label{ex:14-ordena} Orden\textit{á} tu cuarto.
  \ex\label{ex:14-porque} ¿\textit{Por qué no} orden\textit{ás} tu cuarto?
\end{exe}

\subsection{Modalidad epistémica y modalidad deóntica}\label{sec:14-epistemica}

Asimismo, \textcite{palmer1986}, luego de analizar una serie de nociones que
habitualmente se asocian a la de la modalidad, tales como \sig{subjetividad},
\sig{epistémico} / \sig{deóntico}, \sig{factualidad}, \sig{posibilidad} /
\sig{necesidad}, termina inclinándose por distinguir entre modalidad epistémica
y modalidad deóntica, como dos polos básicos que podemos suponer presentes en
todas las lenguas.

\begin{destacado}
La \termino{modalidad epistémica} no solo se aplica a los sistemas modales que
involucran las nociones de \sig{posibilidad} y \sig{necesidad} (como en la
lógica clásica), sino también a todo sistema modal que indique el grado de
compromiso de quien habla con lo que dice, de acuerdo con el nivel de
conocimiento que posee de los acontecimientos a los que se refiere. La
\termino{modalidad deóntica}, por otro lado, es aquella que se asocia a la
expresión de las capacidades, obligaciones, prohibiciones o permisos, por lo
cual incluye dos rasgos: la \sig{subjetividad} (las diversas marcas del «yo» en
aquello que dice) y la \sig{no factualidad}, esto es, la no asertividad: quien
habla no busca transmitir información ni aludir a acontecimientos de la
«realidad», sino que tiene otros propósitos vinculados con lo volitivo.
\end{destacado}

\subsection{Clases de modalidad y tipos de oraciones}\label{sec:14-clases}

De acuerdo con la definición clásica de modalidad como la «actitud subjetiva» del
hablante, se reconoce que las estructuras oracionales codifican, a través de
diferentes recursos formales ---suprasegmentales y morfosintácticos--- y de
ciertos índices léxicos, distintas posiciones enunciativas frente al
\textit{dictum}: el hablante «asevera, conjetura, exhorta, expresa un deseo»
\parencite[98]{kovacci1992}.

La tabla~\ref{tab:14-kovacci} muestra la propuesta de \textcite{kovacci1992},
quien a las dos modalidades distinguidas por \textcite{palmer1986} añade la
modalidad afectiva y se vale de estas tres modalidades para establecer las
diferentes configuraciones oracionales posibles y determinar sus propiedades
formales e índices léxicos característicos.

{\scriptsize\setlength{\tabcolsep}{2pt}
\begin{longtable}{@{}>{\raggedright\arraybackslash}p{23mm}
    >{\raggedright\arraybackslash}p{25mm} >{\raggedright\arraybackslash}p{45mm}
    >{\raggedright\arraybackslash}p{30mm}@{}}
  \caption{Modalidades y configuraciones oracionales \parencite[a partir
    de][]{kovacci1992}}\label{tab:14-kovacci}\\
  \toprule
  \textbf{Modalidades} & \textbf{Configuraciones oracionales} &
    \textbf{Propiedades formales} & \textbf{Índices léxicos} \\
  \midrule
  \endfirsthead
  \toprule
  \textbf{Modalidades} & \textbf{Configuraciones oracionales} &
    \textbf{Propiedades formales} & \textbf{Índices léxicos} \\
  \midrule
  \endhead
  \textbf{Intelectuales o epistémicas}: están relacionadas con el grado de
    conocimiento de los acontecimientos &
    \textbf{Aseverativas o declarativas}: exponen el \textit{dictum} como un
    hecho real &
    Figura tonal que se desarrolla sobre el rango del tono básico y juntura
    terminal descendente. Modo indicativo. &
    Pueden contar con refuerzos de la certeza y convicción del hablante con
    respecto a sus dichos: \textit{por supuesto, ciertamente}\ldots \\
  \addlinespace
  & \textbf{Dubitativas}: ubican el \textit{dictum} en el plano de lo posible,
    probable o hipotético &
    Figura tonal que no es distinta de la declarativa. Pueden incluir formas del
    modo indicativo con interpretación inferencial (futuro, condicional) y del
    modo subjuntivo, además de perífrasis verbales de modalidad (con
    \textit{poder, deber, tener que}). &
    Con índices de actitud: \textit{tal vez, a lo mejor, posiblemente, capaz
    que}\ldots \\
  \addlinespace
  & \textbf{Interrogativas}: indican la actitud de no certeza o de
    desconocimiento con respecto a todo el \textit{dictum} (preguntas totales) o
    a una parte (preguntas parciales) &
    Entonación interrogativa con diferentes configuraciones. Se opone a la
    entonación declarativa por una figura de movimientos alternativos del tono,
    que van desde los agudos a los graves o viceversa. Modo indicativo. &
    Con pronombres interrogativos (\textit{qué, quién}\ldots) y algunos otros
    índices: \textit{¿no?, ¿verdad?} \\
  \midrule
  \textbf{Volitivas}: están asociadas a lo deóntico &
    \textbf{Exhortativas} &
    Presentan variantes tonales e intervalos de altura e intensidad acentual
    según manifiesten un ruego o un mandato categórico. El modo propio de esta
    modalidad es el imperativo. El subjuntivo «completa» las formas vacantes de
    imperativo, aunque también el presente, el futuro del modo indicativo e
    incluso el infinitivo pueden manifestar actitud exhortativa. &
    Diversas interjecciones: \textit{¡vamos!, ¡por favor!, ¡dale!} \\
  \addlinespace
  & \textbf{Desiderativas} &
    Poseen una figura tonal semejante a la declarativa, aunque el hablante puede
    imprimir un particular grado de énfasis. En modo subjuntivo. Con el
    indicativo, solamente la perífrasis \textit{deber / tener que} + infinitivo
    en el pretérito imperfecto o el condicional puede expresar valor
    desiderativo. &
    Con índices desiderativos como \textit{que, ojalá (que)}, en posición
    inicial. \\
  \midrule
  \textbf{Afectivas}: expresan actitudes emocionales (alegría, disgusto,
    sorpresa, reproche\ldots) &
    \textbf{Exclamativas} &
    Las figuras tonales pueden adoptar diferentes configuraciones (ascendentes,
    descendentes, circunflejas) con un refuerzo de intensidad sobre el foco de
    la carga emotiva. Generalmente, se construyen en indicativo. &
    Con pronombres exclamativos, el relativo \textit{lo que}, intensificadores
    como \textit{tan, tanto, tal}\ldots\ y algunas interjecciones. \\
  \bottomrule
\end{longtable}}

Independientemente de la propuesta de clasificación oracional expuesta, es
necesario tener presente que es habitual cierta superposición de matices
modales. Por caso, ante una estructura como \forma{Que pase el que sigue}, es
posible interpretar que existe una intersección entre el deseo (\sig{quiero que
pase el que sigue}) y el mandato (\sig{pido que pase el que sigue}); asimismo,
frente a \forma{¿Cómo pudo haber hecho eso?}, se produce un cruce entre
interrogación y exclamación.

\begin{ejercicio}\label{ej:14-canciones}
Teniendo en cuenta los elementos destacados en cursiva en los fragmentos de
canciones de abajo, identifique la modalidad, el tipo de oración involucrada y
los recursos que la ponen de manifiesto. Dado que los diferentes medios se
integran en la estructuración lingüística, considere tanto los rasgos prosódicos
como la información morfosintáctica y léxica. Especifique cómo contribuye cada
uno de ellos en la interpretación modal predominante de la oración en que
aparecen.

% Letras de canciones protegidas: se incorporan a mano desde el original.
\textoaincorporar{los fragmentos de 1. \textit{Latinoamérica} (Calle 13), con las
formas destacadas a. \textit{se comparte}, b. \textit{que viva}, c. \textit{no
puedes comprar}; 2. \textit{El pianista del gueto de Varsovia} (J.~Drexler), con
a. \textit{tengo}, b. \textit{pude haber sido}, c. \textit{estoy, estabas}; 3.
\textit{Inoportuna} (J.~Drexler), con a. \textit{quien no lo sepa\ldots\
aprenderá}, b. \textit{para, espera, avisa}, c. \textit{tantos, tan}, d.
\textit{eran, decía, Cautela}, e. \textit{cuánta}, f. \textit{tan}}

\begin{resuelto}
1. a. Modalidad deóntica: la predicación \forma{Aquí se comparte}, solo posible
en 3.ª persona, se interpreta \sig{(las cosas) deben ser compartidas}; la
construcción mediopasiva \forma{se comparte} desencadena una lectura de
genericidad. b. Modalidad deóntica, con \forma{que} como índice desiderativo,
que exige el verbo flexionado en modo subjuntivo. c. Modalidad deóntica: la
perífrasis verbal con \forma{poder} en forma negativa indica
\sig{imposibilidad}, y el verbo en 2.ª persona del singular presenta
interpretación genérica (\sig{nadie puede comprar mi vida}).
\end{resuelto}
\end{ejercicio}

\subsection{Una relación compleja: modo y modalidad}\label{sec:14-modo}

Otro punto de confusión en este tema reside en la confluencia entre el
\termino{modo} y la \termino{modalidad}, algo que advierte
\textcite{dahl2015}, también desde un enfoque interlingüístico. El término
\textit{modalidad} alude al amplio dominio que relaciona el contenido
proposicional con la realidad e involucra nociones como \sig{necesidad},
\sig{posibilidad}, \sig{probabilidad}. El \textit{modo}, en cambio, se refiere a
distinciones gramaticales, o sea, a las inflexiones que se incluyen en las
categorías de tiempo, aspecto, modo y evidencialidad, que se expresan en la
sigla \textsc{tame} (véase §\ref{sec:14-evidencialidad}). Estas categorías, en
algunas lenguas como el español, se manifiestan por lo general en los verbos. El
modo, en ese sentido, no siempre resulta de una elección «libre» de quien habla,
sino que suele depender del tipo de oración en que el verbo se inserta; por
ejemplo, la oración desiderativa con \forma{ojalá que} exige, en el español, la
flexión verbal en subjuntivo \exref{ex:14-ojala}:

\begin{exe}
  \ex\label{ex:14-ojala} \textit{Ojalá} hay\textit{a} buen tiempo en el fin de
      semana.
\end{exe}

\textcite{dahl2015} argumenta que, típicamente, en el francés o en el alemán,
por caso, se usa el modo subjuntivo para marcar la \sig{no aserción} de ciertas
oraciones subordinadas: en el francés, aquellas que son complementos de los
predicados de \sig{deseo}; en el alemán, el subjuntivo se emplea, en cambio, en
las oraciones que expresan discurso indirecto.\footnote{Se denomina
\textit{discurso directo} (DD) «el que reproduce de forma literal palabras
pronunciadas, así como pensamientos transcritos en la misma forma en que se
originan» \parencite[§\,43.9a]{rae2025}, y \textit{discurso indirecto} (DI) al
que presenta la voz del hablante que narra lo que él mismo u otras personas han
dicho o pensado; estos contenidos aparecen en subordinadas sustantivas
introducidas por la conjunción subordinante \forma{que}, ausente del DD
\parencite[§\,43.9c]{rae2025}.} Aunque interlingüísticamente las formas
flexivas que no pertenecen al indicativo (considerado el modo «neutro» en el
conjunto de oposiciones que constituye con el subjuntivo y el imperativo en el
español) reflejan una considerable variación, se puede deducir que lo que sí
tienen en común es que se usan en general para expresar la \sig{no aserción}; por
ejemplo, para distinguir lo \sig{posible} de lo \sig{contrafáctico} (contrario a
los hechos efectivamente acaecidos), como en los períodos condicionales, que son
construcciones de dos constituyentes ---una prótasis o condicionante y una
apódosis o condicionado--- que ponen en un vínculo hipotético dos eventos,
\exref{ex:14-parque1} y \exref{ex:14-parque2} (véase el \capref{13},
§\ref{sec:13-condicionales}).

\begin{exe}
  \ex\label{ex:14-parque1} Si \textit{está} soleado, \textit{iré} al parque.
  \ex\label{ex:14-parque2} Si \textit{hubiera estado} soleado, \textit{habría ido}
      al parque.
\end{exe}

Mientras que en \exref{ex:14-parque1} ambos eventos, tanto el \forma{estar
soleado} como el \forma{ir al parque}, se presentan como hechos posibles, el
encadenamiento condicional de \exref{ex:14-parque2}, en cambio, se interpreta
como contrario a los hechos en el pasado: \forma{no estuvo soleado}, por lo
tanto, \forma{no fui al parque}.

\begin{recuadro}[Para ampliar]
\textcite{dahl2015} sostiene que en muchas lenguas emparentadas por razones
genéticas (o sea, lenguas que comparten un «antepasado común»), como el tagalo,
el sundanés o el javanés, de la familia austronesia ---habladas en la isla de
Madagascar, el archipiélago malayo y Oceanía---, se emplean formas de futuro
para comunicar la \sig{irrealidad}. De esto se puede desprender, entonces, la
existencia de un cierto patrón interlingüístico común que opone lo
\textit{realis} frente a lo \textit{irrealis}; y también que este último ámbito
de la manifestación de lo irreal revela una notable variación en lo que atañe a
las funciones específicas que recubre y al repertorio de medios disponibles
según la lengua que analicemos.\pendiente{El original dice «sudanés»}
\end{recuadro}

Ahora bien, ¿qué significa que se gramaticalice la modalidad en una lengua? Así
como algunos verbos plenos pueden ampliar su uso a un comportamiento auxiliar en
las perífrasis verbales, «cediendo» significado léxico y especializándose en
contenido gramatical, también los contenidos gramaticales de los paradigmas
flexivos codifican una información específica de índole muy abstracta. La
regularidad de las oposiciones manifestadas en estos conjuntos cerrados de formas
permite la contraposición interna dentro de un rango posible de valores, que se
expresan necesariamente en determinadas clases de palabras que establecen
conexiones morfosintácticas \parencite{lehmann2002,giammatteo2018}. Por caso, en
el español, la oposición de modo indicativo, subjuntivo e imperativo distingue
tres valores: \forma{cantás / cantes / cantá; comés / comas / comé; vivís / vivas
/ viví}, para la segunda persona voseante del singular de \forma{cantar, comer} y
\forma{vivir}, respectivamente. En esta oposición queda codificado el contraste
\sig{afirmación} / \sig{deseo} / \sig{mandato}, como las diferentes actitudes de
quien habla ante un mismo evento. Forzosamente, estas oposiciones de modo son
exhibidas por los verbos finitos y están vinculadas con las estructuras
sintácticas que integran; por ejemplo, el verbo \forma{comunicar} se combina con
el modo indicativo, en \exref{ex:14-comunica1}, dada la compatibilidad del acto
de «comunicar» con la introducción de una aseveración, que es lo codificado en
la flexión en modo indicativo de \forma{está}; pero rechaza, en este contexto,
el subjuntivo \exref{ex:14-comunica2}, puesto que es incompatible lo que
codifica \forma{esté} (básicamente, no aserción) con el predicado introductor.

\begin{exe}
  \ex\label{ex:14-comunica1} Comunica que est\textit{á} bien.
  \ex[*]{\label{ex:14-comunica2} Comunica que est\textit{é} bien.}
\end{exe}

\begin{ejercicio}\label{ej:14-modos}
A partir de las formas verbales presentes en los siguientes pares de oraciones,
explique qué significados modales se contraponen en ellas, según se muestra en
el ejemplo resuelto~1.

\begin{exe}
  \ex Los novios \textit{son} felices. / \textit{Sean} felices los novios.
  \ex El Señor los \textit{protege}. / El Señor los \textit{proteja}.
  \ex Esta semana \textit{escribiremos} el capítulo. / Esta semana
      \textit{escribiríamos} el capítulo.
  \ex ¡Ojalá \textit{venga} a la fiesta! / ¡Ojalá \textit{viniera} a la fiesta!
  \ex ¡Ojalá \textit{viniera} a la fiesta! / ¡Ojalá \textit{hubiera venido} a la
      fiesta!
  \ex Capaz \textit{juega} mañana. / Capaz \textit{juegue} mañana.
\end{exe}

\begin{resuelto}
1. La oposición de modo indicativo (\forma{son}) y modo subjuntivo
(\forma{sean}) contrapone modalidad epistémica aseverativa frente a modalidad
deóntica de deseo.
\end{resuelto}
\end{ejercicio}

\section{¿Modalidad o evidencialidad?}\label{sec:14-evidencialidad}

Según hemos planteado, la delimitación de la modalidad en una perspectiva
interlingüística no está exenta de controversias, entre otras razones porque la
difusión de la existencia de lenguas que exhiben un \termino{sistema gramatical
evidencial} ha provocado que se abriera el debate sobre si la modalidad y la
evidencialidad constituyen una misma categoría o si son dos categorías
diferentes. Por caso, \textcite{barnes1984} explica que en el tuyuca, una lengua
hablada en Brasil y Colombia, se gramaticalizan los siguientes significados: i)
\textit{evidencia visual}; ii) \textit{evidencia no visual} (proveniente de algún
otro sentido que no sea la vista, predominantemente el oído); iii)
\textit{apariencia} o \textit{deducciones} hechas a partir de la evidencia; iv)
evidencia de segunda mano, es decir, lo que se conoce como el \textit{reporte} o
la \textit{evidencia citativa}, y v) \textit{evidencia inferida}, que se expresa
cuando el hablante posee un conocimiento previo de un estado de cosas o de
ciertos patrones de conducta generales o habituales.

Esto implica que los hablantes del tuyuca, al decir algo, no pueden no marcar
con un afijo gramatical ---así como nosotros en español lo hacemos con el tiempo
o con el modo--- cuál ha sido la fuente de acceso a su conocimiento de los
acontecimientos referidos: si esa información proviene de algo que han visto
«con sus propios ojos», que han oído, que deducen de ciertos indicios, que les
han contado o que infieren a partir de otros datos.

Como vemos, entonces, hay lenguas de distintas familias (entre ellas, amerindias,
eslavas como el macedonio, túrquicas\ldots) que expresan algunos de estos
contenidos mediante categorías morfológicas específicas.\pendiente{El original
dice «amerindias, macedónicas, turcas»} Según \textcite{willett1988}, lo más
recurrente es que se distinga, en la morfología verbal, si la fuente para la
evidencia es directa o «atestiguada», indirecta (o citativa) o inferida (es
decir, la que se deriva de asignar causas, en función de una construcción
mental, a las situaciones observadas). Para estas lenguas, las marcas de
evidencia atestiguada proceden ---mediante fenómenos de fusión y reducción
sintáctica y fonológica--- de verbos de percepción sensorial; las de reporte, de
verbos de habla; y las de inferencia, de tiempos perfectos o de verbos de
observación visual. Este hecho muestra que las marcas evidenciales reflejan un
proceso que se conoce como \termino{gramaticalización}
\parencite{company2016}, por el que cierto material léxico (como los verbos de
\sig{percepción} y \sig{comunicación}) se vuelve, mediante un procedimiento de
reanálisis, un afijo gramatical. Una forma de gramaticalización que guarda
semejanzas con lo detectado por Willett se puede comprobar también en el
español: el verbo \forma{haber} en \forma{amar he > amaré, amar has >
amarás\ldots} se transformó, por un cambio que se produce a lo largo de la
historia de la lengua, en el sufijo de flexión del tiempo futuro; es decir,
\forma{haber}, que en un primer estadio «auxilia» al infinitivo en \forma{amar
he} (= \sig{he de amar}), termina integrándose en el verbo léxico, como parte de
la estructura morfológica de palabra.

En el caso de las marcas evidenciales, es importante destacar que, aunque sean
parte de la flexión de los verbos, su alcance es oracional: su ámbito de
operación recae sobre toda la estructura sintáctica, tal como ocurre, en el
español, con la flexión de tiempo, que se expresa en el verbo, pero «temporaliza»
toda la predicación. En efecto, la marcación evidencial especifica la fuente de
la información del hablante para la aserción y, por tanto, concierne a toda la
oración aseverada, aunque solamente se manifieste en un afijo gramatical.

Hemos visto, hasta aquí, que la evidencialidad, concebida como una categoría
lingüística (no meramente gramatical), involucra la fuente de información
(directa o indirecta ---inferencial y reportativa o citativa---), mientras que
la modalidad epistémica se refiere a los significados modales que manifiestan un
soporte, un respaldo completo, parcial o neutro (en aquellos casos en los que la
proposición no es objeto de una evaluación epistémica) de la aserción. Por
consiguiente, aunque se puede reconocer que modalidad y evidencialidad son
categorías conceptualmente diferentes, existe una interacción funcional entre
ambas: la epistemicidad \parencite{bybee2012}.

\subsection{Evidencialidad, deixis y perspectivización}\label{sec:14-deixis}

Por otra parte, en la codificación lingüística de la evidencialidad, no podemos
omitir la importancia de la \termino{deixis}. \textcite[319]{lyons1997} trata la
deixis como una operación de señalamiento referencial, es decir, una indicación
de «la relación que se establece entre las expresiones lingüísticas y lo que
representan en el mundo (o universo del discurso) en ocasiones particulares de
enunciado». En sintonía con esto, \textcite{bermudez2016} defiende que la
evidencialidad supone coordenadas que son propias de la deixis espacial
---\textit{puntos de referencia}, \textit{dirección} y \textit{distancia}---, que
se reinterpretan, metafóricamente, en el dominio evidencial:

\begin{itemize}
  \item los \textit{puntos de referencia} oscilan en un continuo entre el
        hablante, como fuente única, y una fuente general que condensa el
        conocimiento compartido por una comunidad, es decir, su
        \textit{folklore};
  \item la \textit{dirección} es el modo de acceso al conocimiento, que puede ser
        \sig{hacia} el hablante, a través de impresión sensorial, o \sig{desde}
        el hablante, mediante inferencia o razonamiento;
  \item la \textit{distancia} concierne, en un polo, a la información recibida de
        primera mano y, en el polo opuesto, a lo que proviene del rumor, o sea,
        de una fuente poco fidedigna.
\end{itemize}

\begin{destacado}
En función de esta confluencia entre la deixis y el sistema de perspectiva, la
evidencialidad puede ser entendida, según \textcite{hassler2010}, como un
\termino{fenómeno deíctico de perspectivización} desde el punto de vista del
hablante, que el oyente debe reinterpretar desde su propia posición discursiva.
\end{destacado}

Esta perspectivización implícita en la deixis permite distinguir, para la
autora, modalidad epistémica de evidencialidad: la primera aporta
monológicamente la actitud de quien habla; la otra requiere del oyente la
elaboración de la referencia a la fuente de información manifestada por el
hablante y la realización de todos los ajustes necesarios sobre la fiabilidad de
las fuentes.

En síntesis, el reconocimiento de estos sistemas evidenciales ha llevado a
\textcite{palmer1986} a admitir que es frecuente que no pueda establecerse una
distinción clara entre \textit{juicios} (asociados tradicionalmente a la noción
de modalidad en la lógica) y \textit{evidenciales} (marcas gramaticales del
acceso a la información), dado que los juicios de quien habla están naturalmente
atados al grado de fiabilidad de la evidencia que posee.

\section{Evidencialidad: paradigma frente a estrategias}\label{sec:14-estrategias}

En cuanto al español, podemos asumir, como plantea
\textcite{gonzalezvazquez2016}, que la discusión no queda zanjada, porque se
presentan las siguientes dificultades:

\begin{itemize}
  \item persiste la controversia en cuanto a qué relación guarda la modalidad con
        la evidencialidad;
  \item se admite que tanto en los sistemas gramaticalizados como en los no
        gramaticalizados los marcadores evidenciales pueden denotar
        polisémicamente varios significados;
  \item se usaron, a lo largo de la historia, diferentes criterios para
        clasificar los evidenciales, aunque la clasificación más extendida sea la
        de \textcite{willett1988}, que presentamos en
        §\ref{sec:14-evidencialidad}.
\end{itemize}

El recorrido que sigue expone la dificultad de estudiar la modalidad sin atender
a cómo inciden otras categorías asociadas a ella en la flexión verbal, como el
tiempo y el aspecto, y sin tener en cuenta la evidencialidad, categoría que, en
las lenguas románicas, no se organiza estrictamente en torno a oposiciones
flexivas y, en cambio, se expresa mediante una amplia variedad de mecanismos.
Como el español no es una lengua que cuente con un paradigma evidencial
sistematizado en las formas verbales, se considera que pertenece al grupo de
lenguas que desarrollan diversas \termino{estrategias evidenciales}, mediante el
empleo de ciertos ítems léxicos (verbos como \forma{oír, ver, parecer} y
adverbios como \forma{claramente, visiblemente, aparentemente}, entre otros) o
recursos gramaticales diversos, tales como ciertos usos del futuro, del
condicional y del pretérito imperfecto, o la configuración oracional de la
impersonalidad, entre todos los que ilustramos a continuación.

\subsection{El pretérito imperfecto}\label{sec:14-imperfecto}

\textcite[24]{hassler2017}, por ejemplo, analiza el pretérito imperfecto en
textos narrativos y periodísticos como marcador evidencial de que la información
procede de otra fuente, distinta de quien la reproduce, a partir de ejemplos
como \exref{ex:14-papa} y \exref{ex:14-wall}. De esta manera, el locutor se
distancia de sus enunciados y realiza simultáneamente una atenuación del
contenido proposicional:

\begin{exe}
  \ex\label{ex:14-papa} El Papa Francisco \textit{recibía} el lunes en el
      Palacio Apostólico a la última tanda de prelados de las provincias
      eclesiásticas del sur de España, incluidas las Islas Canarias, Ceuta y
      Melilla. (GlossaNet, \textit{ABC}, 13/03/2014).
  \ex\label{ex:14-wall} Wall Street abrió hoy sin tendencia definida y el Dow
      Jones de Industriales \textit{bajaba} un 0,07~\%, con los mercados
      cautelosos a la espera del inicio de la temporada de resultados
      trimestrales. (GlossaNet, \textit{El País}, 10/04/2014).
\end{exe}

Este uso del pretérito imperfecto ha sido relacionado con su función citativa
solo dentro del marco de sus valores pragmáticos y discursivos en el discurso
oral \parencite{reyes1990}; por ejemplo, cuando formulamos una pregunta como
\forma{¿A qué hora hab\textbf{ía} que salir?}, que busca confirmación de una
información ya dicha. Para \textcite{hassler2017}, a partir de una intersección
semántica entre la atenuación y la evidencialidad, el pretérito imperfecto ya no
expresa sus significados primarios o funciones prototípicas
\parencite[véase][]{giammatteo2018trombetta}, sino otros significados
secundarios que están relacionados con la actitud epistémica del hablante frente
a su enunciado y a la fuente de información, lo que se puede subsumir bajo la
noción de \termino{evidencialidad indirecta}.

Este significado secundario del pretérito imperfecto puede derivarse de su
aspecto imperfectivo, por medio del cual se contempla una acción en su
transcurso y desarrollo, lo que hace que la validez de la oración se presente
también como abierta o indeterminada: la aspectualidad imperfectiva conduce a la
modalidad debido a su indeterminación. Sin embargo, en otras lenguas, como el
alemán, donde no hay formas imperfectivas, esta función es asumida por verbos
modales que, además de su función deóntica, tienen usos evidenciales (por
ejemplo, el verbo \forma{wollen}).

Por otra parte, con una clara interpretación atenuativa de modalidad de
\sig{irrealidad}, el pretérito imperfecto puede emplearse en todas las lenguas
románicas como recurso de la cortesía lingüística \parencite[21]{hassler2017},
según los ejemplos reproducidos como \exref{ex:14-fr}, \exref{ex:14-port} y
\exref{ex:14-it}:

\begin{exe}
  \ex\label{ex:14-fr} fr. Je \textit{voulais} vous demander, Monsieur, de me
      permettre d'assister à votre cours.\\
      \sig{[Yo] quería pedirle, señor, me permita asistir a su clase}.
  \ex\label{ex:14-port} port. \textit{Podias} passar este texto a limpo?\\
      \sig{¿Podías pasar este texto en limpio?}
  \ex\label{ex:14-it} it. Scusi, \textit{volevo} un bicchier d'acqua.\\
      \sig{Disculpe, quería un vaso de agua}.
\end{exe}

Como muestran los datos de arriba y las glosas que explicitan el significado,
también en español es un recurso cortés el empleo del imperfecto como mitigador
de la fuerza del mandato. Las formas verbales \forma{quería} y \forma{podía}
indican cortesía en la expresión de lo pedido: \forma{asistir a su clase},
\forma{pasar este texto en limpio} y \forma{pedir un vaso de agua}, en
\exref{ex:14-fr}, \exref{ex:14-port} y \exref{ex:14-it}, respectivamente.

\subsection{El futuro y el condicional}\label{sec:14-futuro}

También \textcite[22]{hassler2017} muestra que el futuro simple puede usarse con
efectos de atenuación, sobre todo en contextos de cortesía, para suavizar una
petición, como en \exref{ex:14-dira}:

\begin{exe}
  \ex\label{ex:14-dira} Usted me di\textit{rá} lo que tengo que hacer. (En vez
      de \textit{Dígame lo que tengo que hacer}).
\end{exe}

En otros contextos, el uso del futuro simple puede mitigar el grado de certeza
con que se expresa un contenido proposicional, como en los empleos consabidos de
las formas de futuro con interpretación temporal de presente
\exref{ex:14-chico}, que sugieren que la aserción se asienta sobre inferencias:

\begin{exe}
  \ex\label{ex:14-chico} Ese chico ahora tend\textit{rá} unos diez años.
\end{exe}

\textcite[30]{hassler2017} se concentra en el \termino{futuro narrativo}
retrospectivo, al que pone en relación con la evidencialidad. Se trata de un uso
infrecuente en español, que aparece para referirse a situaciones pasadas con
respecto al momento de habla, pero que son presupuestas,\footnote{En el análisis
semántico, la presuposición alude al hecho de que hay expresiones lingüísticas
incrustadas dentro de una matriz que no se ve afectada por la negación: así,
tanto \forma{Me interesa que se estén jugando las Olimpiadas} como \forma{No me
interesa que se estén jugando las Olimpiadas} presentan como presupuesto el
hecho de que efectivamente se están realizando los juegos olímpicos. Dicho de
otro modo, las subordinadas introducidas por predicados que expresan reacciones
emotivas (\forma{gustar, interesar, alegrarse\ldots}) no quedan bajo el ámbito
de la negación, y su contenido, que se da por conocido y cuyo valor de verdad no
es cuestionado, es presentado «bajo el prisma de una evaluación, una emoción
[\ldots]» \parencite[§\,25.3z]{rae2025}.} que se exponen como dadas por el
locutor y que, como no son deícticas, son incompatibles con complementos
adverbiales referidos al momento presente \exref{ex:14-ragazzo}.\pendiente{Comentario
de F. Carranza: precisar la relación entre el futuro narrativo, la
evidencialidad y la presuposición}

\begin{exe}
  \ex\label{ex:14-ragazzo} it. Il ragazzo partì all'improvviso. \textit{Tornerà}
      solo dopo tre anni / *fa tre anni.\\
      \sig{El muchacho partió imprevistamente. Regresará solo después de tres
      años / *hace tres años}.
\end{exe}

Con este empleo del futuro narrativo retrospectivo, en general en tercera
persona, quien habla «simula» una dislocación retrospectiva con respecto al
momento de habla. En este sentido, Hassler interpreta que se puede atribuir al
futuro narrativo un valor evidencial: el hablante se presenta en tales
enunciados como alguien que conoce los hechos y comunica sus conclusiones, porque
conoce la historia hasta su final.

Como estamos viendo, la morfología verbal flexiva del español queda disponible
para sumar nuevas lecturas, no solamente las de modalidad, como los usos de
cortesía, sino también las que, más recientemente, tienden a ser analizadas como
empleos evidenciales \parencite{bermudez2016,garcianegroni2016,reyes1990}, como
el futuro inferencial \exref{ex:14-nadie}; el futuro citativo, que remite a algo
ya dicho y adopta valor concesivo \exref{ex:14-fortuna}, y el futuro citativo
refutativo, que se construye con la forma perifrástica
\exref{ex:14-conductor}:\pendiente{El original cita «García Negroni (2016)»;
la entrada es García Negroni y Marcovecchio (2016)}

\begin{exe}
  \ex\label{ex:14-nadie} No hab\textit{rá} nadie todavía, porque la puerta está
      cerrada. (= \sig{Supongo que no habrá\ldots}).
  \ex\label{ex:14-fortuna} Hab\textit{rá} trabajado mucho, pero esa fortuna no se
      hace con el trabajo de un año. (Réplica a una intervención previa, que el
      hablante concede parcialmente).
  \ex\label{ex:14-conductor} A: ---Es un excelente conductor.\\
      B: ---¿Qué \textit{va a ser} buen conductor? Es un irresponsable.
      (Réplica).
\end{exe}

De manera semejante al futuro, también admiten estos usos el condicional y el
imperfecto, cuando se emplean para hacer una suposición (valor inferencial)
\exref{ex:14-impacto} o retomar algo dicho (valor citativo)
\exref{ex:14-llamabas}:

\begin{exe}
  \ex\label{ex:14-impacto} El impacto se hab\textit{ría} producido por el estado
      de embriaguez del que conducía. (Como parte de la información de una
      crónica periodística).
  \ex\label{ex:14-llamabas} ¿Cómo te llama\textit{bas}? Me lo dijiste, ya sé,
      pero olvidé tu nombre.
\end{exe}

\begin{destacado}
De los datos anteriores se desprende la polifuncionalidad de muchas de las formas
verbales del español y, en lo relativo a la evidencialidad, la «ganancia» de
interpretaciones secundarias que sugieren, fundamentalmente, un acceso a la
información que se comunica mediante razonamiento deductivo ---por ejemplo, en
\exref{ex:14-nadie}, \forma{la puerta cerrada} funciona como un indicio para
desencadenar la inferencia \forma{no habrá nadie todavía}--- o a través de la
mención de dichos ajenos.
\end{destacado}

\begin{ejercicio}\label{ej:14-formas}
Identifique las formas verbales destacadas de las siguientes oraciones y
explique, en términos modales y evidenciales, las oposiciones que plantean,
según se muestra en el ejemplo resuelto~1.

\begin{exe}
  \ex \textit{Tenía} sesenta años cuando lo conocí. / \textit{Tendría} sesenta
      años cuando lo conocí.
  \ex Juan no está en la casa: ¿\textit{ha ido} a la facultad esta noche? /
      ¿\textit{habrá ido} a la facultad esta noche?
  \ex ---No encuentro mis llaves.\\
      ---Las \textit{dejaste} / \textit{habrás dejado} en la puerta.
  \ex El gobierno \textit{llamará} / \textit{llamaría} a elecciones
      extraordinarias.
  \ex Si lo \textit{hubiera sabido} / \textit{sabía}, no \textit{habría venido} /
      \textit{venía} a la reunión.
  \ex ¡Buen día! \textit{Quiero} / \textit{Quisiera} saber si hay alguna chance de
      probarla antes de comprar. Tengo una 6º y no me \textit{resulta} /
      \textit{estaría resultando}. (Consulta sobre una boquilla de un
      instrumento musical en un blog).\pendiente{El original dice «una 6*»; el
      asterisco se reemplazó por «º» para evitar que se lea como juicio de
      gramaticalidad: verificar}
\end{exe}

\begin{resuelto}
1. Las formas verbales \forma{tenía} (imperfecto) y \forma{tendría}
(condicional) oponen \sig{conocimiento efectivo} frente a \sig{conocimiento
inferido}.
\end{resuelto}
\end{ejercicio}

\subsection{Los auxiliares de perífrasis}\label{sec:14-auxiliares}

Algunos de los auxiliares de las perífrasis verbales poseen empleos modales y
evidenciales. Como se vio en el \capref{11}, \forma{poder, deber
(de)}\footnote{\forma{Deber} con \forma{de} o sin \forma{de} se revela, en el
uso del español general, como un fenómeno afectado por la variación. En lo
normativo, se recomienda \forma{deber} para el empleo que indica necesidad u
obligación, y \forma{deber de} para la expresión de probabilidad, pero el uso no
siempre registra esta diferencia.} y \forma{tener que} son habitualmente
reconocidos como índices de la modalidad deóntica, que manifiesta
\sig{habilidad}, \sig{permiso}, \sig{necesidad} u \sig{obligación}, y también de
la modalidad epistémica, que comprende desde lo \sig{posible} hasta lo
\sig{probable}. Pero, además, dichos auxiliares también pueden analizarse como
manifestación del grado de compromiso que el hablante asume según la fiabilidad
de los indicios que le permiten realizar ciertas deducciones
\parencite{cornillie2007}.

\begin{exe}
  \ex\label{ex:14-billetera} La billetera \textit{puede / debe de / tiene que}
      haber quedado en el auto.
\end{exe}

Los diferentes auxiliares señalan la menor o mayor fuerza de las presunciones de
quien habla: con \forma{poder} se sugiere que los indicios son débiles para una
aseveración, mientras que con \forma{deber} y \forma{tener que} se indica que
las «pistas» sobre dónde quedó la billetera permiten una deducción con mayor
grado de probabilidad. También con un auxiliar como \forma{soler} se produce una
imbricación entre el contenido aspectual de \sig{frecuencia} y el de modalidad
epistémica \parencite{giammatteo2011}.

\begin{exe}
  \ex\label{ex:14-suele} El profesor \textit{suele} venir por las tardes, pero hoy
      no concurrió al instituto.
\end{exe}

Asimismo, un predicado de \sig{comunicación} como \forma{prometer}, según vimos
en el \capref{11}, muestra un empleo auxiliarizado que sugiere un cierto grado de
adhesión de quien habla con respecto a la proyección de un acontecimiento
futuro:

\begin{exe}
  \ex\label{ex:14-promete} El encuentro \textit{promete} ser una fiesta.
\end{exe}

Nuevamente, así como planteamos con respecto a la morfología verbal, comprobamos
en algunos auxiliares una cierta hibridación entre contenidos modales y otros,
como el aspecto \exref{ex:14-suele}, o la descripción de la fuerza ilocutiva de un
acto de habla \exref{ex:14-senalo2}.\pendiente{El original remite aquí a (33) y
(34), que ilustran otra cosa} Ciertamente, la morfología verbal y las
perífrasis verbales constituyen dos extremos de la expresión de contenidos
gramaticales: en un caso, se recurre a la afijación; en otro, a la combinación de
dos palabras en una misma unidad verbal. Pero el comportamiento de ambos
mecanismos es complementario \parencite{marcovecchio2010}:

\begin{exe}
  \ex\label{ex:14-senalo1} El profesor le señaló que complet\textit{ara} el
      ejercicio.
  \ex\label{ex:14-senalo2} El profesor le señaló que \textit{debía} completar el
      ejercicio.
\end{exe}

Tanto en \exref{ex:14-senalo1} como en \exref{ex:14-senalo2}, la oración
dependiente de \forma{señalar} expresa la fuerza deóntica, pero en
\exref{ex:14-senalo1} lo hace mediante el uso del subjuntivo \forma{completara} y
en \exref{ex:14-senalo2} recurre al auxiliar \forma{debía} en modo indicativo, lo
que muestra la conexión supletoria entre la flexión en modo subjuntivo y el verbo
auxiliar modal.

\begin{ejercicio}\label{ej:14-parafrasis}
Establezca paráfrasis del significado de las siguientes oraciones utilizando
perífrasis verbales, según se muestra en el ejemplo resuelto~1.

\begin{exe}
  \ex El ACV \textit{es prevenible} y \textit{tratable}.
  \ex Eso \textit{es indiscutible}.
  \ex Las emociones \textit{no se medican}: \textit{se procesan}.
  \ex Esta prenda \textit{no es lavable} en máquina.
  \ex \textit{No se vive} del aire, sin trabajar.
  \ex En esta sala, \textit{no se grita}.
\end{exe}

\begin{resuelto}
1. El ACV \textit{puede ser prevenido y tratado} / \textit{se puede prevenir y
tratar}.
\end{resuelto}
\end{ejercicio}

\begin{ejercicio}\label{ej:14-gramaticalizacion}
Explique, en términos de gramaticalización, la distribución del modo subjuntivo
de la forma simple destacada o de la forma perifrástica en indicativo.

\begin{exe}
  \ex Dados los resultados del laboratorio, el médico le indicó al paciente que
      \textit{hiciera} (*\textit{hacía}) una dieta.
  \ex Dados los resultados del laboratorio, el médico le indicó al paciente que
      \textit{tenía que hacer} (*\textit{tuviera que hacer}) una
      dieta.\pendiente{¿*\textit{tuviera que hacer}? Parece aceptable:
      revisar el juicio}
\end{exe}
\end{ejercicio}

\subsection{Adverbios, locuciones y oraciones subordinadas}\label{sec:14-adverbios}

Ahora bien, el español, como las demás lenguas románicas, también se vale de
adverbios o locuciones equivalentes, así como de ciertas oraciones subordinadas
que, ubicados en la periferia oracional, se comportan como «estrategias» para
codificar contenidos evidenciales \parencite{marcovecchio2001}. En lo formal,
estos constituyentes quedan delimitados por un contorno entonativo propio,
indicio de su alcance oracional, extrapredicativo, y suelen ubicarse en el margen
preverbal, aunque justamente la delimitación melódica propia les brinda gran
movilidad posicional. En lo interpretativo, se trata de recursos que, por un
desplazamiento metafórico desde contenidos léxicos o menos gramaticalizados
hacia otros que manifiestan significados gramaticales o más gramaticalizados,
dejan de funcionar en el ámbito de la predicación y pasan a aludir a la esfera
del hablante.

Así, \forma{por ahí} es un SP que señala un contenido espacial en
\exref{ex:14-porahi1}, pero adquiere otra interpretación, vinculada a la
modalidad y a las inferencias que realiza quien enuncia, en
\exref{ex:14-porahi2}. Las huellas de la gramaticalización también se advierten
en la pronunciación, ya que en la interpretación modal ambos elementos pasan a
emitirse como una unidad y con desplazamiento acentual.

\begin{exe}
  \ex\label{ex:14-porahi1} Las llaves están \textit{por ahí}. (Pron.
      \fonol{poɾ aˈi}).
  \ex\label{ex:14-porahi2} \textit{Por ahí} (≡ \textit{capaz, a lo mejor}) viene
      esta tarde a visitarme. (Pron. \fonol{poɾ aˈi} \textasciitilde\
      \fonol{poɾˈai}).
\end{exe}

Justamente, acudiendo a la noción de evidencialidad, se advierte que la locución
\forma{por ahí} puede manifestar tanto el carácter citativo de un enunciado (a
cuyo contenido, habitualmente, se contrapone la opinión de quien habla)
\exref{ex:14-misa} como la lectura hipotético-inferencial que se le imprime a un
enunciado \exref{ex:14-calle} \parencite{marcovecchio2015}:\pendiente{El
original cita «Marcovecchio, 2015, p. 2020»}

\begin{exe}
  \ex\label{ex:14-misa} «\textit{Por ahí se dice} que no tendría que haber ido [a
      la misa]. Pero yo fui invitado y me parece que era una buena
      oportunidad» para homenajear a Juan Pablo II, uno de los principales
      responsables de que la Argentina y Chile no hayan entrado en guerra hace 30
      años. (lanacion.com.ar, 24/12/2008).
  \ex\label{ex:14-calle}
    \begin{xlist}
      \ex\label{ex:14-calle-a} Para finalizar, González advirtió: «No hay delito,
          ella está declarando porque desapareció. Si ella estuvo durmiendo en la
          calle, \textit{por ahí} no vio televisión [\ldots]».
          (elintransigente.com, 10/07/2018).
      \ex\label{ex:14-calle-b} (Estoy en condiciones de deducir que) Dado que
          durmió en la calle, es probable que no haya visto televisión.
    \end{xlist}
\end{exe}

Además, en \exref{ex:14-calle-a}, en que \forma{por ahí} manifiesta un vínculo
inferencial entre \forma{haber estado durmiendo en la calle} y \forma{no ver
televisión}, se puede llegar a establecer una correspondencia con las oraciones
subordinadas llamadas \textit{causales de la enunciación}, que tienen valor modal
en tanto se vuelven recursos epistémico-evidenciales \exref{ex:14-calle-b}
\parencite{marcovecchio2001}.

Se trata de un caso de evidencialidad inferida (el modo de acceso es el
razonamiento, no la experiencia sensible), y no es casual que en la locución
\forma{por ahí} sea un elemento deíctico el factor crucial para disparar la
lectura evidencial. El acercamiento a la esfera personal y subjetiva produce un
debilitamiento asertivo que es propio de procesos que convierten ciertos
elementos, en principio concernientes a la descripción de un estado de cosas en
el mundo, en recursos gramaticales asociados a la expresión de significados
relativos a la situación de habla y de los interlocutores: lo que se conoce como
\termino{pragmaticalización}.

Con respecto a las oraciones subordinadas (estudiadas en el \capref{13}), su
empleo como marcadores de la modalidad o de la evidencialidad no constituye una
excepción, sino más bien lo contrario: los períodos oracionales constituidos por
prótasis y apódosis ---todos los relacionados con el ámbito de la
\sig{causalidad}, que incluye subordinadas causales, consecutivas, finales,
condicionales y concesivas (véase el \capref{13},
§\ref{sec:13-logicas})--- son favorables a desencadenar lecturas que conciernen a
las inferencias que realiza el hablante \exref{ex:14-delpo} o a su propia
situación como «diciente» \exref{ex:14-murray}
\parencite{kovacci1992,marcovecchio2016}:

\begin{exe}
  \ex\label{ex:14-delpo}
    \begin{xlist}
      \ex Delpo se la jugó. Y jugó un partidazo. Le tocó sufrirlo a Ferrer, que
          también tenía sus inconvenientes ---un esguince en el tobillo
          izquierdo, por el que se infiltró para jugar---, pero no por eso era un
          adversario menos peligroso; a fin de cuentas, sus buenas razones tenía
          el alicantino \textit{para haber llegado hasta aquí también}.
          (\textit{La Nación}, Buenos Aires, 04/07/2013).
      \ex \sig{Si llegó hasta aquí también, (será porque / puedo deducir que)
          sus buenas razones tenía el alicantino}.
    \end{xlist}
  \ex\label{ex:14-murray}
    \begin{xlist}
      \ex «\textit{Para ser sincero}, ni siquiera yo lo sé», dijo Murray.
          (Canchallena, \textit{La Nación}, Buenos Aires, 21/01/2014).
      \ex \sig{Si hablo con sinceridad, (te digo que) ni siquiera yo lo sé}.
    \end{xlist}
\end{exe}

\begin{ejercicio}\label{ej:14-subjuntivo}
Explique qué significados modales manifestados en el modo verbal (e incluso en
las conjunciones) se contraponen en los siguientes pares de oraciones (por
ejemplo, \sig{conocimiento de primera mano} frente a \sig{conocimiento
inferido}, \sig{aserción} frente a \sig{deseo}, etc.), según se muestra en el
ejemplo resuelto~1.

\begin{exe}
  \ex Insiste en que los chicos siempre \textit{se portan} bien. / Insiste en que
      los chicos siempre \textit{se porten} bien.
  \ex Siento que \textit{cierra} esa puerta. / Siento que \textit{cierre} esa
      puerta.
  \ex Creo que no \textit{asistirá} a la clase. / No creo que \textit{asista} a
      la clase.
  \ex El juez no cree que \textit{existe} un delito. / El juez no cree que
      \textit{exista} un delito.
  \ex Depende de la historia que \textit{tengo} entre manos. / Depende de la
      historia que \textit{tenga} entre manos.
  \ex Literariamente, el fútbol es muy fértil y está bueno siempre que nos
      \textit{conduce} a otros lugares. / Literariamente, el fútbol es muy
      fértil y está bueno siempre que nos \textit{conduzca} a otros lugares.
  \ex \textit{Si} no \textit{presentás} el pasaporte, no salís del país. /
      \textit{Como} no \textit{presentes} el pasaporte, no salís del país.
\end{exe}

\begin{resuelto}
1. Las formas verbales en oposición \forma{portan} (indicativo) / \forma{porten}
(subjuntivo) contraponen modalidad epistémica aseverativa a modalidad deóntica
exhortativa.
\end{resuelto}
\end{ejercicio}

\subsection{Marcadores discursivos}\label{sec:14-marcadores}

Si bien el término \textit{marcadores discursivos} (concepto introducido en el
\capref{12}, §\ref{sec:12-marcadores}) se emplea habitualmente en el análisis
gramatical, no alude a una clase de palabras (como las categorías nombre o verbo,
por caso), sino a un conjunto categorialmente heterogéneo de ítems (adverbios,
conjunciones, locuciones prepositivas, incluso verbos, todos con otras funciones
vinculadas con la interacción) que comparten entre sí propiedades discursivas.
Estos atributos pueden sintetizarse en la «aislabilidad» respecto de la
sintaxis predicativa y la no dependencia del significado léxico que aporta la
predicación sobre la que inciden; el impacto interaccional sobre los intercambios
comunicativos, por el que se evoca a los participantes de la situación
dialógica; y su comportamiento ilocutivo, en tanto expresan el «estado mental»
de quienes hablan en lo relativo a sus emociones, el grado de conocimiento de los
acontecimientos referidos, el acuerdo o el desacuerdo. En síntesis, se
relacionan directamente con el acto de habla \parencite{sanchezlopez2020}.

Dentro del amplio grupo de marcadores discursivos, \forma{evidentemente, por lo
visto, al parecer} ilustran un subgrupo particular: los \termino{evidenciales
citativos e inferenciales}, que \textcite[8--9]{cornillie2015}, por ejemplo,
estudian en el español peninsular. Consisten en una serie de marcadores del
discurso referidos a la percepción visual, pero que se emplean para indicar una
inferencia o también una información de oídas. El análisis de estos autores
sigue un criterio discursivo interaccional y muestra que existe una asimetría
entre \forma{evidentemente}, por un lado, y \forma{por lo visto} y \forma{al
parecer}, por otro: mientras que el primero se asocia al señalamiento de
pensamientos compartidos \exref{ex:14-vidente}, los otros dos expresan una
evidencia indirecta, \exref{ex:14-coca} y \exref{ex:14-nazis}.

\begin{exe}
  \ex\label{ex:14-vidente} Entonces, yo lo que hice fue ponerme en contacto con
      la madre primero\ldots\ que ya la\ldots\ había habla<(d)>o con ella en otras
      ocasiones, entonces, \textit{evidentemente} no le conté lo de la muerte,
      pero yo pensé que si ella llamaba a la policía judicial la harían más caso
      que si llamaba yo, de parte de una vidente de las muchísimas que han
      llama<(d)>o, porque han llama<(d)>o\ldots\ mogollón [\ldots].
  \ex\label{ex:14-coca} <H2> Sabes que ayer pillaron a un\ldots\ un director de
      un banco con\ldots\\
      <H3> Sí.\\
      <H2> \ldots\ kilo y medio de coca, ¿no?\\
      <H1> Con quince kilos.\\
      <H4> Quince, quince [kilos de cocaína]. Sí. 15, 15.\\
      <H2> ¡Ah, quince! Ah. Se me\ldots\ se me había corrido la coma. <todos> ah,
      15. [\ldots]\\
      <H2> Está bien, está bien quince, ¿no?\\
      <H4> Está muy bien. Era un distribuidor de primera línea, \textit{por lo
      visto}.
  \ex\label{ex:14-nazis} <H3> De todas maneras, lo que sí que\ldots\ ha ido a más
      en Europa eh\ldots\ han sido los brotes de xenofobia\ldots\\
      <H1> Aquí no han llega<(d)>o aún, ¿no?\\
      <H4> Sí, pero Madrid se iba a convertir el domingo, o se va a convertir, ya
      veremos, en una\ldots\ van a venir todos los nazis de Europa, \textit{por
      lo visto}, o neonazis y esto va a ser\ldots\ un festejo nazi. Por eso,
      \textit{al parecer}, queda prohibido\ldots
\end{exe}

\pendiente{Falta la fuente de los ejemplos orales (corpus citado por Cornillie y
Gras Manzano)}

Los datos de \exref{ex:14-coca} y \exref{ex:14-nazis} sugieren que \forma{ser
distribuidor de primera línea}, por un lado, y \forma{venir todos los nazis de
Europa} y \forma{quedar prohibido el festejo}, por otro, bajo la incidencia de
\forma{por lo visto} y \forma{al parecer}, no comunican una percepción visual de
los hechos, sino una deducción que realiza quien habla a partir del conocimiento
de acontecimientos que son la base del razonamiento.\pendiente{El original dice
«por lo visto y evidentemente»} En \exref{ex:14-coca}, que la carga de cocaína
decomisada a un director de un banco sea de quince kilos (y no de un kilo y
medio) es la causa para suponer que se trataba de un gran distribuidor; en
\exref{ex:14-nazis}, que los brotes de xenofobia en Europa se hayan incrementado
desencadena la suposición de que un encuentro de manifestantes previsto en
Madrid se vuelva un festejo neonazi y que, por lo tanto, las autoridades lo
prohíban.

Como podemos ir advirtiendo, el uso de los marcadores discursivos está modelado
por razones dialectales y sociolectales (o sea, reflejan un potencial de
variación de acuerdo con la habitualidad de su empleo en diferentes variedades
del español según su distribución geográfica o según el grado de formalidad de la
situación comunicativa). Por esto, dado también que el contenido interaccional de
esos marcadores es abstracto, muchas de estas unidades no son conmutables en el
uso concreto de comunidades diversas de hablantes del español. Por ejemplo, si
comparamos el empleo de \forma{igual, lo mismo} y \forma{justo} en el español
peninsular y en el argentino
\parencite{garcianegroni2013,marcovecchio2022,marcovecchio2023}, encontramos que
los usos discursivos peninsulares pueden resultar opacos o no «naturales» en
nuestra variedad, más allá de otros valores compartidos. Por ejemplo,
\forma{justo}, en ambas variedades, puede explicitar una reformulación de
equivalencia que establece el hablante para los dos segmentos en que se
interpone \exref{ex:14-justo1}:

\begin{exe}
  \ex\label{ex:14-justo1} Simplemente, las pone en su sitio, \textit{justo} (≡
      \textit{precisamente, exactamente}) en su verdadera dimensión.
\end{exe}

Sin embargo, otros datos muestran que no siempre nos resultan «transparentes»
\forma{igual, lo mismo} y \forma{justo}: en \exref{ex:14-igual} y
\exref{ex:14-lomismo}, \forma{igual} y \forma{lo mismo} se comportan como
marcadores epistémicos asemejables al empleo de \forma{capaz (que)} en el Río de
la Plata; y en \exref{ex:14-justo2}, \forma{justo} expresa, en una interacción
dialógica, una confirmación de lo dicho previamente por otro, es decir, la
adhesión con el interlocutor:

\begin{exe}
  \ex\label{ex:14-igual} Lleva paraguas, que \textit{igual} llueve. (= \textit{Capaz
      que} llueve).
  \ex\label{ex:14-lomismo} ---¿Dónde está el profesor?\\
      ---No sé. \textit{Lo mismo} está en su despacho. (= \textit{Capaz que} está
      en su despacho).
  \ex\label{ex:14-justo2} ---¿Estás insinuando que el director va a dimitir?\\
      ---\textit{Justo}. (= \textit{Precisamente}).
\end{exe}

El ejemplo \exref{ex:14-justo2} exhibe un uso reformulador de \forma{justo}. De
hecho, la \termino{reformulación}, proceso que permite que el hablante vuelva
sobre un discurso anterior para presentarlo desde una nueva perspectiva o para
reinterpretarlo, es un campo fértil para considerar una extensión de funciones
interaccionales de ítems como \forma{es decir, o sea, a saber, mejor dicho} o
\forma{en / con otras palabras}, que claramente son expresiones que explicitan
esta operación metalingüística de establecimiento de lazos de relativa
equivalencia entre dos segmentos. No obstante, \textcite{rodriguezramalle2014},
que estudia los reformuladores \forma{es decir} y \forma{o sea} frente a
\forma{esto es, a saber, mejor dicho, con otras palabras}, etc., concluye que la
distribución de unos y otros en el uso no es semejante. Según la autora,
\forma{es decir} y \forma{o sea} se distinguen de los demás por poder ser
utilizados por un hablante distinto del que realiza la primera intervención y,
asimismo, pueden introducir conclusiones, como se ve en \exref{ex:14-ingeniero}
y \exref{ex:14-analogia}, fragmentos de intercambios dialógicos en los que los
reformuladores introducen la voz de un hablante distinto con respecto a los
dichos previos, que motivan una reinterpretación que se sugiere como resultado
del establecimiento de inferencias generalizadoras y conclusivas:

\begin{exe}
  \ex\label{ex:14-ingeniero} A: Cuando yo era niño, mi madre escuchaba música
      occidental [\ldots]. Luego mi padre compró un fonógrafo eléctrico y el
      sonido me pareció totalmente distinto\ldots\ Aquella diferencia hizo que me
      decidiese a hacerme ingeniero electrónico.\\
      B: ¿\textit{O sea que} se hizo ingeniero electrónico a causa de la música?
  \ex\label{ex:14-analogia} A: Que los niños son analogistas y sobrepasan siempre
      el grado de regularidad del sistema de la lengua adulta lo sabe todo el
      mundo.\\
      B: \textit{Es decir que} es la facultad de analogía sobre todo lo que
      consideras como innato. (Casado, 1991, pp.~107--108).\pendiente{Casado
      (1991) no figura en la bibliografía}
\end{exe}

También \textcite{rodriguezramalle2014} analiza como marcadores evidenciales las
locuciones conjuntivas \forma{conque, de manera que} y \forma{así que}, en
relación con los reformuladores \forma{es decir} y \forma{o sea}, dado que también
pueden introducir un discurso que remite a las fuentes de información (sin
especificaciones sobre si se trata de evidencialidad citativa o inferencial) y
que motiva la emisión de los dichos que siguen. De ahí las interpretaciones
expresadas en \exref{ex:14-conque-b} y \exref{ex:14-alvaro-b}:

\begin{exe}
  \ex\label{ex:14-conque}
    \begin{xlist}
      \ex ¡\textit{Conque} por fin te vas y nos dejas! $\rightarrow$ ¿\textit{Así
          que / o sea que} por fin te vas y nos dejas?
      \ex\label{ex:14-conque-b} \sig{Me he enterado, me han dicho que por fin te
          vas y nos dejas}.
    \end{xlist}
  \ex\label{ex:14-alvaro}
    \begin{xlist}
      \ex \textit{Es decir que} tú eres el famoso Álvaro, del que todo el mundo
          habla. $\rightarrow$ \textit{De manera que} tú eres el famoso Álvaro,
          del que todo el mundo habla.
      \ex\label{ex:14-alvaro-b} \sig{Me han dicho o he deducido por lo que me han
          contado que tú eres el famoso Álvaro}.
    \end{xlist}
\end{exe}

\begin{destacado}
Como vemos, es amplio el repertorio de recursos disponibles para explicitar o
sugerir significados modales y evidenciales. Además, el análisis de estos
contenidos gramaticales requiere atender a la variación de uso, de acuerdo con
diferencias geográficas, situacionales y de registro, así como a la extensión de
funciones y significados de ítems que, por procesos de pragmaticalización, suman
comportamientos discursivos.
\end{destacado}

\begin{ejercicio}\label{ej:14-evidenciales}
Indique qué tipo de evidencialidad y / o de modalidad codifican las formas
destacadas en las siguientes oraciones, según se muestra en el ejemplo
resuelto~1.

\begin{exe}
  \ex \textit{En principio}, no \textit{sería} nada grave [\ldots].
  \ex \textit{Debe ser} un fenómeno multicausal. La lógica \textit{indicaría} en
      principio que, al crecer la pobreza, \textit{deberían aumentar} este tipo
      de modalidades. \textit{Al parecer} lo único que está creciendo es el robo
      tipo motochorro.
  \ex ¿Vos te \textit{llamabas}\ldots?
  \ex ---¡\textit{Qué va a ser} hermoso! Es horrible.
  \ex El señor \textit{sabe venir} por aquí a la mañana.
  \ex Otro abogado explicó: «\textit{Puede que} los minoristas no \textit{acepten}
      ahora, pero \textit{si se concretara} el canje, algunos \textit{podrían
      aceptarlo}». (\textit{La Nación}, 16/04/2013).
\end{exe}

\begin{resuelto}
1. Tanto la locución \forma{en principio} como la flexión en condicional de
\forma{ser} codifican evidencialidad citativa y / o inferencial (la
interpretación es ambigua).
\end{resuelto}
\end{ejercicio}

\begin{destacado}
\textbf{Para tener en cuenta.} Estas diferentes actitudes subjetivas se combinan
con la \termino{polaridad} oracional: afirmativa o negativa. La polaridad
afirmativa no necesita una marca explícita; en cambio, la negativa, que siempre
expresa un distanciamiento, un cuestionamiento por parte de quien habla con
respecto a lo que dice, requiere marcas específicas (\forma{no, ni, nunca, nada,
nadie\ldots}), los llamados \termino{términos de polaridad negativa}.
\end{destacado}

\section{Modalidad del propio acto de enunciación}\label{sec:14-enunciacion}

Al estudiar la estructura oracional, podemos advertir que, típicamente, hay
constituyentes periféricos con respecto a la predicación que expresan contenidos
relativos a la situación de habla, lo que concierne tanto a los interlocutores
como a la misma estructuración del discurso. En el primer caso, se trata de
expresiones como \forma{francamente, sinceramente, en confianza, para ser
honesto\ldots} En el segundo, nos referimos a todos los elementos que establecen
indicios de la presencia del hablante como la persona que organiza los
contenidos en la secuencia discursiva: ordenándolos y jerarquizándolos
(\forma{primero / en primer lugar, segundo / en segundo lugar, por / para
empezar, en conclusión / para concluir, en síntesis\ldots}) o estableciendo
reformulaciones ---de equivalencia o no--- entre segmentos (\forma{en otras
palabras, dicho de otro modo, por poner un ejemplo, para ilustrar\ldots}).

\begin{ejercicio}\label{ej:14-tipos}
Reconozca las propiedades formales y los ítems léxicos que permiten incluir las
oraciones de abajo dentro de los tres tipos de modalidad diferenciados, según se
muestra en el ejemplo resuelto~1.

\begin{exe}
  \ex Modalidades intelectuales o epistémicas:
    \begin{xlist}
      \ex Por supuesto, guardé las llaves.
      \ex Por ahí, dejé las llaves sobre el escritorio.
    \end{xlist}
  \ex Modalidades volitivas (deónticas): exhortativa, desiderativa:
    \begin{xlist}
      \ex Que pase el que sigue.
      \ex Chicos, (les pido que) guarden los útiles en sus mochilas.
      \ex Chicos, pueden / deben guardar los útiles en sus mochilas.
      \ex Chicos, ¡a ver si guardan los útiles en sus mochilas (de una vez por
          todas)!
      \ex ¡Ojalá fuera rica! Si fuera rica, viajaría por el mundo.
    \end{xlist}
  \ex Modalidades afectivas: exclamativas:
    \begin{xlist}
      \ex ¡Qué lindo (que) es soñar!
      \ex Siempre dice qué (/ lo) lindo (que) es soñar.
    \end{xlist}
\end{exe}

\begin{resuelto}
1. a) Modo indicativo, con refuerzo de la aserción \forma{por supuesto}. b) Modo
indicativo, con marcador hipotético inferencial \forma{por ahí}.
\end{resuelto}
\end{ejercicio}

\begin{ejercicio}\label{ej:14-periferia}
Explique si los elementos destacados dependen del verbo de la predicación
oracional o se orientan hacia el acto de enunciación. Justifique su respuesta.

\begin{textolectura}
El hecho habría ocurrido a escasos kms. de San Guillermo. [\ldots]
\textit{Extraoficialmente}, se presume que el robo rondaría en una cifra cercana
a 300 mil pesos. [\ldots] \textit{Según fuentes consultadas por Radio Belgrano},
habría sido abordado por dos sujetos [\ldots]. Los delincuentes se habrían
llevado la recaudación.
\end{textolectura}
\end{ejercicio}

\begin{recuadro}[Reflexión final]
Sin lugar a duda, las sutilezas de la expresión de la modalidad requieren
continuar ahondando en los posibles significados que se explicitan o que apenas
se sugieren. No obstante, hasta aquí hemos confrontado datos provenientes del
uso real, que ilustran los recursos prosódicos, morfosintácticos, léxicos y de
la configuración oracional que inciden en la manifestación de la modalidad y de
la evidencialidad, como «estrategias» que permiten la reinterpretación de
ciertos morfemas de la flexión verbal, de algunos adverbios y conjunciones del
español en términos de empleos citativos o inferenciales. En este recorrido,
hemos consolidado temas tratados en capítulos anteriores (las categorías
morfológicas, las perífrasis verbales, la estructura oracional con funciones
nucleares ---argumentales o no, pero siempre dentro del ámbito del SV--- y
periféricas, las alternancias cuasirreflejas mediopasivas e impersonales) y los
hemos reconsiderado integrados en la combinatoria oracional. Asimismo, hemos
destacado la necesidad de atender al potencial de variación que existe en el
español: los empleos reales de la lengua muestran las particularidades de uso de
los recursos gramaticales asociados a la modalidad y la evidencialidad.
\end{recuadro}

% --- Secciones finales -------------------------------------------------
\lecturasrecomendadas

\solucionario

\begin{solucion}{\ref{ej:14-canciones}}
2. a. Modalidad epistémica aseverativa con verbos flexionados en presente del
modo indicativo. b. Modalidad epistémica dubitativa: la perífrasis verbal con
\forma{poder} encadenada a un infinitivo compuesto indica un hecho contrafáctico
(véase §\ref{sec:14-modo}) o lo que no sucedió (\sig{no fui el pianista del gueto
de Varsovia}). c. Modalidad epistémica aseverativa, con verbos en presente y
pretérito imperfecto del modo indicativo.

3. a. Modalidad deóntica exhortativa articulada en la oración principal sobre el
verbo \forma{aprender} en futuro, con valor de verdad general; la subordinada
introducida por el relativo \forma{quien} presenta el verbo en modo subjuntivo,
\forma{sepa}, que adopta valor de \sig{inespecificidad}, porque alude a
cualquier persona que no esté en conocimiento del contenido referido. b.
Modalidad epistémica aseverativa negativa con sentido general, con el verbo en
presente de indicativo. c. Modalidad afectiva, con intensificador resaltado por
la repetición (\forma{tantos, tan}). d. Modalidad epistémica declarativa con
verbos en indicativo (\forma{eran, decía}); \forma{Cautela} manifiesta modalidad
deóntica exhortativa con valor de advertencia. e. Modalidad afectiva, con el
intensificador \forma{cuánta}. f. Modalidad afectiva, con uso explícito y
enfático del pronombre de 2.ª persona y empleo del intensificador \forma{tan}
repetido.\pendiente{Las letras de los apartados (a, b, c, e, f, g en el
enunciado original; a--f en el solucionario) no coinciden}
\end{solucion}

\begin{solucion}{\ref{ej:14-modos}}
2. La oposición de modo indicativo (\forma{protege}) y modo subjuntivo
(\forma{proteja}) opone modalidad epistémica aseverativa a modalidad deóntica
desiderativa. 3. La oposición de futuro (\forma{escribiremos}) y condicional
(\forma{escribiríamos}), ambos dentro del modo indicativo, opone modalidad
epistémica aseverativa frente a modalidad epistémica dubitativa. 4. La oposición
de presente de subjuntivo (\forma{venga}) e imperfecto de subjuntivo
(\forma{viniera}) opone modalidad deóntica de deseo posible frente a modalidad
deóntica de deseo imposible (en el presente: \forma{¡Ojalá viniera hoy!}) o
improbable (en el futuro: \forma{¡Ojalá viniera uno de estos días!}). 5. La
oposición de imperfecto (\forma{viniera}) y pluscuamperfecto de subjuntivo
(\forma{hubiera venido}) contrapone modalidad deóntica de deseo imposible (en el
presente) o improbable (en el futuro) a modalidad de deseo irreal
contrafáctico. 6. La oposición de presente de indicativo (\forma{juega}) y
presente de subjuntivo (\forma{juegue}) contrapone modalidad epistémica de
incerteza débil frente a modalidad epistémica de incerteza fuerte.
\end{solucion}

\begin{solucion}{\ref{ej:14-formas}}
2. Las formas verbales \forma{ha ido} (perfecto compuesto) y \forma{habrá ido}
(futuro perfecto) oponen una pregunta que se formula desde el
\sig{desconocimiento} a otra que ofrece una inferencia. 3. Las formas verbales
\forma{dejaste} (perfecto simple) y \forma{habrás dejado} (futuro perfecto)
oponen \sig{conocimiento efectivo} a \sig{conocimiento inferido}. 4. Las formas
verbales \forma{llamará} (futuro) y \forma{llamaría} (condicional) oponen
\sig{conocimiento efectivo} a \sig{evidencialidad citativa} o bien
\sig{conocimiento inferido}. 5. Las formas verbales \forma{hubiera sabido / no
habría venido} (pluscuamperfecto de subjuntivo y condicional compuesto) y
\forma{sabía / venía} (ambas en imperfecto de indicativo) oponen registro formal
a registro coloquial informal; en ambos pares se mantiene la interpretación
contrafáctica. 6. Las formas \forma{quiero / resulta} (presente de indicativo) y
\forma{quisiera / estaría resultando} (imperfecto de subjuntivo y condicional)
oponen grados de cortesía: \forma{quisiera} y \forma{estaría resultando}, al
indicar \sig{irrealidad}, atenúan la fuerza deóntica del pedido o el comentario
negativo sobre la boquilla.
\end{solucion}

\begin{solucion}{\ref{ej:14-parafrasis}}
2. Eso \textit{no puede / no debe ser discutido} / \textit{no se puede
discutir}. 4. Esta prenda \textit{no puede / no debe ser lavada} / \textit{no se
puede / debe lavar} en máquina. 5. No \textit{se puede vivir} del aire, sin
trabajar. 6. En esta sala, \textit{no se puede / no se debe
gritar}.\pendiente{El original numera estas respuestas 2 a 5 y no resuelve
la oración 3 (\textit{Las emociones no se medican})}
\end{solucion}

\begin{solucion}{\ref{ej:14-gramaticalizacion}}
En las lenguas, se reconoce una especie de «complementariedad» entre recursos
gramaticales, aquí morfológicos (la flexión en modo subjuntivo de
\forma{hiciera}) o morfosintácticos (la perífrasis verbal \forma{tenía que
hacer}). El verbo \forma{indicar} introduce un COD que expresa una recomendación
mediante la modalidad deóntica; por lo tanto, el verbo de la subordinada o bien
manifiesta esa modalidad con el verbo flexionado en subjuntivo, o bien lo hace
con la perífrasis verbal de modalidad deóntica.
\end{solucion}

\begin{solucion}{\ref{ej:14-subjuntivo}}
2. En combinación con \forma{cierra} (indicativo), \forma{sentir} es un verbo de
\sig{percepción sensorial} que manifiesta modalidad epistémica aseverativa; con
\forma{cierre} (subjuntivo), \forma{sentir} se convierte en verbo psíquico y se
interpreta como \sig{lamentar}, lo que manifiesta modalidad afectiva. 3. Las
formas verbales \forma{asistirá} (indicativo) y \forma{asista} (subjuntivo)
contraponen modalidad epistémica aseverativa de certeza débil a modalidad
epistémica dubitativa. 4. Las formas verbales \forma{existe} (indicativo) y
\forma{exista} (subjuntivo) contraponen modalidad epistémica aseverativa con
posicionamiento fuerte del hablante (\sig{el juez no cree, pero el hablante
afirma que existe delito}) a un posicionamiento neutro de quien habla (\sig{el
juez no cree en la existencia del delito, y el hablante no asume tampoco un
compromiso epistémico}). 5. Las formas verbales \forma{tengo} (indicativo) y
\forma{tenga} (subjuntivo) contraponen modalidad epistémica aseverativa acerca de
la existencia de la historia mencionada al desconocimiento de su posible
existencia o la falta de especificidad al mencionarla. 6.\pendiente{El original
no resuelve la oración 6 (\textit{siempre que nos conduce / conduzca})} 7.
Ambas prótasis condicionales en oposición expresan un vínculo hipotético entre
\forma{la presentación del pasaporte} y \forma{la posibilidad de salir del país};
no obstante, con \forma{si} más predicado en modo indicativo, el posicionamiento
del hablante es neutro, mientras que con \forma{como} más predicado en modo
subjuntivo se sugiere un posicionamiento de advertencia o amenaza.
\end{solucion}

\begin{solucion}{\ref{ej:14-evidenciales}}
2. Las perífrasis con \forma{deber}, la flexión en condicional de
\forma{indicar} y \forma{deber} y el marcador discursivo \forma{al parecer}
expresan una intersección entre modalidad epistémica de probabilidad y
evidencialidad inferencial. 3. La forma en pretérito imperfecto del modo
indicativo \forma{llamabas} sugiere evidencialidad citativa. 4. La forma
perifrástica \forma{ir a} + infinitivo sugiere evidencialidad citativa, a la vez
que la no adhesión a los dichos previos. 5. La perífrasis verbal con
\forma{saber} expresa aspecto habitual y también involucra el matiz modal de
posibilidad. 6. \forma{Puede que} introduce una predicación modalizada
epistémicamente, con indicación de \sig{posibilidad}, que exige la forma
flexionada de \forma{aceptar} en modo subjuntivo. El período condicional que
sigue, con la prótasis en pretérito imperfecto del modo subjuntivo y la apódosis
en condicional, indica \sig{irrealidad}, específicamente, algo improbable en el
futuro.
\end{solucion}

\begin{solucion}{\ref{ej:14-tipos}}
2. Modalidades volitivas (deónticas): exhortativa, desiderativa. a) Modo
subjuntivo exigido por el índice de exhortación o deseo \forma{que}. b) Modo
subjuntivo exigido por el verbo \forma{pedir}, que explicita la fuerza
ilocutiva; inclusión del vocativo para nombrar a las personas a las que va
dirigida la exhortación. c) Perífrasis verbales que expresan modalidad deóntica
y empleo del vocativo para nombrar a las personas a las que va dirigida la
exhortación. d) Uso de estructuras que pueden expresar modalidad deóntica
(\forma{a ver si\ldots}), más un marcador del énfasis del pedido (\forma{de una
vez por todas}); empleo del vocativo para nombrar a las personas a las que va
dirigida la exhortación. e) Modo subjuntivo exigido por el índice desiderativo.

3. Modalidades afectivas: exclamativas. a) Uso de pronombres enfáticos que se
corresponden con la curva entonativa exclamativa. b) Uso de pronombres enfáticos
que indican la modalidad exclamativa.
\end{solucion}

\begin{solucion}{\ref{ej:14-periferia}}
\forma{Extraoficialmente} y \forma{según fuentes consultadas por Radio Belgrano}
no dependen de los verbos de las oraciones sobre las que inciden porque no
admiten focalización ni establecen alguna función sintáctica dentro de la
predicación. Desempeñan funciones periféricas y expresan que las declaraciones
hechas provienen de información obtenida de segunda mano.
\end{solucion}
